% Part: first-order-logic
% Chapter: proof-systems
% Section: tableaux

\documentclass[../../../include/open-logic-section]{subfiles}

\begin{document}

\iftag{FOL}
      {\olfileid[fr]{fol}{prf}{tab}}
      {\olfileid[fr]{pl}{prf}{tab}}

\olsection{\usetoken{P}{tableau}}

Alors que de nombreux systèmes de !!{derivation} manipulent des
agencements d'!!{sentence}s, les !!{tableau}s manipulent des
!!{signed formula}s. Une !!{signed formula} est un couple formé d'un
signe de valeur de vérité ($\True$ ou $\False$) et d'un !!{sentence} :
\[
\sFmla{\True}{!A} \text{ ou } \sFmla{\False}{!A}.
\]
Un !!{tableau} est constitué de !!{signed formula}s disposées dans un
arbre qui se ramifie vers le bas. Il commence par un certain nombre
d'\emph{hypothèses}, puis se poursuit par des !!{signed formula}s
obtenues à partir de l'une des !!{signed formula}s situées au-dessus,
par application d'une règle d'inférence. Chaque règle permet d'ajouter
une ou plusieurs !!{signed formula}s à l'extrémité d'une branche,
ou deux !!{signed formula}s côte à côte. Dans ce dernier cas, la branche
se divise en deux, et les deux formules ajoutées constituent les
extrémités des deux nouvelles branches.

Les règles de décomposition propositionnelles, appliquées à une
!!{signed formula} complexe, ajoutent des !!{signed formula}s dont les
formules sont des sous-!!{formula}s immédiates. Elles vont par paires,
une règle pour chacun des deux signes. Par exemple, la règle
$\TRule{\True}{\land}$ s'applique à $\sFmla{\True}{!A \land !B}$ et
permet d'ajouter les deux !!{signed formula}s $\sFmla{\True}{!A}$
et~$\sFmla{\True}{!B}$ à l'extrémité de toute branche contenant
$\sFmla{\True}{!A \land !B}$. La règle $\TRule{\False}{\land}$
permet quant à elle de diviser une branche en ajoutant côte à côte
$\sFmla{\False}{!A}$ et $\sFmla{\False}{!B}$.
Un !!{tableau} est fermé si chacune de ses branches contient une paire
de !!{signed formula}s $\sFmla{\True}{!A}$ et $\sFmla{\False}{!A}$
portant sur la même formule.
\begin{editorial}
Note de l'édition française : la description par sous-formules est
limitée ici aux règles de décomposition propositionnelles. La coupure,
présentée plus loin dans l'original, n'est pas une telle règle.
Le nom de la règle de conjonction fausse est corrigé : l'original
y place toute la formule au lieu du seul connecteur. Dans le tableau
ci-dessous, les deux dernières étapes portent également le nom de la
règle de conjonction vraie, et non celui de l'implication vraie.
\end{editorial}

La relation $\Proves$ fondée sur les !!{tableau}s se définit ainsi :
$\Gamma \Proves !A$ si et seulement s'il existe un ensemble fini
$\Gamma_0 = \{!B_1, \dots, !B_n\} \subseteq \Gamma$ pour lequel il
existe un !!{tableau} fermé ayant pour hypothèses
\[
\{\sFmla{\False}{!A}, \sFmla{\True}{!B_1}, \dots, \sFmla{\True}{!B_n}\}.
\]
Voici, par exemple, un !!{tableau} fermé qui établit que
$\Proves (!A \land !B) \lif !A$ :
\begin{oltableau}
  [\sFmla{\False}{(\formula{A} \land \formula{B}) \lif \formula{A}}, just = \TAss
    [\sFmla{\True}{\formula{A} \land \formula{B}}, just = {\TRule{\False}{\lif}[1]}
      [\sFmla{\False}{\formula{A}}, just = {\TRule{\False}{\lif}[1]}
        [\sFmla{\True}{\formula{A}}, just = {\TRule{\True}{\land}[2]}
          [\sFmla{\True}{\formula{B}}, just = {\TRule{\True}{\land}[2]}, close]
        ]
      ]
    ]
  ]
\end{oltableau}

Un ensemble $\Gamma$ est incohérent dans le calcul des !!{tableau}s
s'il existe un !!{tableau} fermé ayant pour hypothèses
\[
\{\sFmla{\True}{!B_1}, \dots, \sFmla{\True}{!B_n}\}
\]
pour certains $!B_i \in \Gamma$.

Les !!{tableau}s ont été inventés indépendamment dans les années 1950
par Evert Beth et Jaakko Hintikka, puis simplifiés et popularisés par
Raymond Smullyan. Ils sont très faciles à utiliser, car leur construction
suit une procédure très systématique. Cette propriété les rend aussi
faciles à mettre en œuvre sur ordinateur. Toutefois, un !!{tableau}
est souvent difficile à lire, et son rapport avec une démonstration
n'est pas toujours facile à discerner. La méthode est par ailleurs très
générale : de nombreuses logiques différentes possèdent des systèmes de
!!{tableau}s. Ceux-ci aident aussi à trouver des !!{structure}s qui
satisfont des !!{sentence}s donnés, ou des ensembles d'!!{sentence}s.
Si l'ensemble est satisfaisable, il n'admet aucun !!{tableau} fermé :
tout !!{tableau} possède donc une branche ouverte. On peut « lire »
sur une branche ouverte une !!{structure} qui satisfait l'ensemble,
à condition que toutes les règles susceptibles d'être appliquées sur
cette branche l'aient été. Le lien avec le calcul des séquents est
également très étroit : pour l'essentiel, un !!{tableau} fermé est une
!!{derivation} condensée du calcul des séquents, écrite dans l'autre sens.

\end{document}
